\documentclass[12pt]{article}
\usepackage[utf8]{inputenc}

\usepackage{../mystyle}

\title{Exceptional Collection Note}
\author{Joseph Sullivan}
\date{November 2025}

\DeclareMathOperator{\Kom}{Kom}
\DeclareMathOperator{\Cyl}{Cyl}
\DeclareMathOperator{\Stab}{Stab}
\DeclareMathOperator{\Slice}{Slice}
\DeclareMathOperator{\chern}{ch}
\DeclareMathOperator{\todd}{td}

\DeclareMathOperator{\Kos}{Kos}

\DeclareMathOperator{\ext}{ext}
%\DeclareMathOperator{\hom}{hom}

\newcommand{\cQ}{\mathcal{Q}}
\newcommand{\cZ}{\mathcal{Z}}
\newcommand{\LCom}{\mathcal{LC}}

\newcommand{\rddots}{\text{\reflectbox{$\ddots$}}}



\begin{document}

\maketitle

\begin{prop}
    Let $X\hookrightarrow \P^n_\C$ be a smooth closed subscheme of dimension $m<n$. Then, $\cO_X(-n),\dots,\cO_X$ can never be an exceptional collection.
\end{prop}

\begin{proof}
    I want to exhibit non-vanishing for some
    \[\Ext^i(\cO_X, \cO_X(-k)) = H^i(X,\cO_X(-k))\]
    for any $i$ and $1\leq k \leq n$. It seems to me that the least likely cohomology to vanish is $H^m(X,\cO_X(-n))$, which by Serre duality is $H^0(X, \cO_X(n) \otimes \omega_X)$. We argue that $\omega_X$ is $n$-regular, by showing that
    \[H^i(X,\omega_X(n-i)) = 0,\]
    for $i>0$, which follows from Kodaira vanishing (since we only consider up to $i=m$, and $n>m$). So, $\omega_X(n)$ is globally generated, so being nonzero as a sheaf means that
    \[H^0(X,\cO_X(n)\otimes \omega_X) \neq 0.\]
    Thus, the collection is not exceptional.
\end{proof}


\end{document}